\section{Representative Samples are Optimal in Large-Budget Pilots}
\label{sec:large-budget}

We begin our analysis of optimal pilot design by considering the optimal pilot sample composition for large pilot budgets.
We show that as the pilot budget grows $b\to\infty$, the set of normalized optimal pilot designs converges to a representative sample---one that matches the composition of the target population $\bm{s}_1 \propto \bm{s}_3$, for any fixed follow-up sample $\bm{s}_2$.
Intuitively, a large pilot allows treatments to continue if and only if they have positive effects in the pilot sample, so we should want them to have positive effects in the pilot if and only if they do in the target population.
This is ensured uniquely by a large representative sample.

We first state our main result for the large-budget regime: as the pilot budget $b$ increases without bound, the investment of optimal pilots grows without bound ($I\to\infty$) and the normalized composition of any optimal pilot converges to the normalized representative design $\bm{s}_3 / \bm{c}^\top \bm{s}_3$.
\begin{theorem}[Large-budget Limiting Optimal
  Policy]
  \label{thm:large-budget-limiting-optimal-policy}
  Suppose the prior distribution of $\bm\tau$ has a density positive on an open ball containing the origin, and the average treatment effect on the target population has finite second moment under the prior.
  Then in the large-budget limit $b\to\infty$, the investment of optimal pilots diverges
  \[
    \inf_{\bm s_1^\star\in\mathcal{S}_1^{\star}(b)} \bm
    c^\top \bm s_1^\star
    \to \infty
  \]
  and the unit-investment normalized optimal designs converge to the normalized representative design,
  \[
    \sup_{\bm s_1^\star\in\mathcal{S}_1^{\star}(b)}
    \left\|
    \frac{\bm s_1^\star}{\bm c^\top \bm
    s_1^\star}-\frac{\bm s_3}{\bm c^\top \bm s_3}
    \right\|
    \to 0.
  \]
\end{theorem}

See \Cref{\sectionof{thm:large-budget-limiting-optimal-policy-appendix}} for the
\hyperref[thm:large-budget-limiting-optimal-policy-appendix]{proof}.
We prove the result by showing a more general large-budget result for objectives of the form
\[
  V_Z(\bm s_1)\defeq \E[\Test_1 \cdot Z\cdot \ATE_3],
\]
where $Z\geq 0$ is a bounded weight that is non-adaptive (independent of the pilot outcome conditional on the true treatment effect), has positive expectation conditional on the true treatment effect $\E[Z\mid\bm\tau]>0$ almost surely.
The pilot impact objective is the special case $Z=\Test_2$, while the impact of adopting based on a single RCT is the special case $Z=1$.
Therefore, representative sampling is the unique large-budget limiting impact-optimal design whether or not there is a follow-up RCT before adoption.

The proof proceeds in two stages.
First, we show that the oracle performance achieved by continuing exactly when the target population effect $\mathrm{ATE}_3$ is positive can be matched by a noise-free ``infinite-budget'' pilot if and only if the pilot sample is representative $\bm{s}_1 \propto \bm{s}_3$, since only then does the pilot not make any type I or type II errors, as visualized in \Cref{fig:large-budget-visualization}.
Second, we show that as the pilot investment grows, its performance is uniformly approximated by the performance of the noise-free pilot.
Performance is bounded away from the oracle performance both for any fixed non-representative pilot sample composition and if the pilot investment is capped, so the optimal pilot converges to a representative sample with ever-increasing investment.
This result and the idea of the noise-free pilot match one intuitive reason why representative samples might be desirable in pilot studies: we should want exactly those treatments that have positive effects in the target population to receive a follow-up study.
We can only approximate this rule when the pilot budget is large.
This sets the stage for a setting where representative pilots are not optimal.

\vspace{0.2cm}
\begin{figure}[ht]
  \centering
  \includegraphics[width=0.8\textwidth]{figures/large_budget_visualization_left.pdf}
  \caption{
    \textbf{Why a representative sample is optimal for large budgets.}
    When the sample is non-representative, i.e. $\bm{s}_1 \npropto \bm{s}_3$, the pilot decision boundary is misaligned with the boundary separating positive and negative average treatment effects in the target population.
    Even with an infinite sampling investment, there are conditional average treatment effects $\bm{\tau}$ for which the pilot incorrectly classifies the sign of the target population average treatment effect $\mathrm{ATE}_3$.
  }
  \label{fig:large-budget-visualization}
\end{figure}
